Consider any given communication round. 
We omit the communication round $n$ from the notation to simplify the notation. 
We define the size $K$ sequence of loss functions denoted as $\{ F_i^{(j)} \}_{j = 0}^{K - 1}$, where $K$ is the number of local update steps.
Recall that $F_i^{(k)}$ is the local objective function of the $i$-th user evaluated on data from the $k$-th local update step.
Let $g_i^{(k)} = \nabla F_i^{(k)} \left( \theta_i^{(k)} \right)$ denote the gradient of user $i$ on the $k$-th local update step.
Equation~\ref{eq: theta} defines the standard local update equation for the FedAvg ($K > 1$) or FedSGD ($K = 1$) procedure\edit{, which are produced by users performing stochastic gradient descent for $K$ updates}. 
\begin{align}
    \theta_i^{(k)} = \theta_i^{(k - 1)} - \eta g_i^{(k - 1)} \textrm{ for all } k = 1, \ldots, K
    \label{eq: theta}
\end{align}

After $K$-steps of local training, we aggregate the updated models $\theta_i^{(K)}$ or, equivalently, the changes made by local updates as in Equation~\ref{eq: global}.
\edit{The new shared model becomes the averaged element-wise parameters from the locally trained user models.}
\begin{align}
        \theta_g \leftarrow \dfrac{1}{\abs{S}} \sum_{i \in S} \theta_i^{(K)} = \theta_g + \dfrac{1}{\abs{S}} \sum_{i \in S} \underbrace{\left( \theta_i^{(K)} - \theta_g \right)}_{-\eta\, \times\, g_{Method}}
    \label{eq: global}
\end{align}

Our comparison of various methods will focus on understanding their expressions of the $\theta_i^{(K)} - \theta_g$ term given in Equation~\ref{eq: global}.
\edit{This expression is the change that a user's local training made to the shared model.}
For the FedAvg procedure, we can use Equation~\ref{eq: theta} recursively to write 
\begin{align}
    \theta_i^{(K)} = \theta_i^{(0)} - \eta \sum_{j = 0}^{K - 1} g_i^{(j)}  \label{eq: recursive}
\end{align}
Noting $\theta_i^{(0)} = \theta_g$, the change a user makes with the FedAvg procedure, $\theta_i^{(K)} - \theta_g$, is driven by $\sum_{j = 0}^{K - 1} g_i^j$.
We refer to this component of a ususer'shange as the \textbf{method gradient}, $g_{Method}$, due to its formulation and use in Equation~\ref{eq: recursive}.
Similarly, other FL methods have a method gradient consisting of the gradients at local update steps.
Equation~\ref{eq: methods} presents the specific gradients used from  Equation~\ref{eq: theta} for FedSGD, FedAvg \citep{fedavg}, and FOMAML \citep{reptile}.
\begin{align}
    g_{FedSGD} = g_i^{(0)}, \quad g_{FedAvg} = \sum_{j = 0}^{K - 1} g_i^{(j)}, \quad g_{FOMAML} = g_i^{(K - 1)} \label{eq: methods}
\end{align}

\edit{
Notably, the above methods result in local changes to the global model based on sums of different local gradients computed to update the model. 
In the next section, we demonstrate which and how much of each local gradient used during local training determines how much emphasis is placed between minimizing Equation~\ref{eq: ims} and Equation~\ref{eq: rp}. 
Furthermore, we discover that within-round learning rates can flexibly emphasize objectives of initial model success and rapid personalization based on the decay used. 
In Section~\ref{sec: decay-exp}, we show that choosing decay that places additional emphasis on Equation~\ref{eq: ims}, relative to FedAvg and FOMAML, results in solutions with improved personalized performance, despite Equation~\ref{eq: rp} being the objective associated with rapid personalization.  
}